\section{Related Work}
\label{sec: related_work}
\paragraph{Reinforcement Learning with Verifiable Rewards.}  Recent large-scale reasoning models such as OpenAI o1~\citep{openai2024learning}, DeepSeek‑R1~\citep{guo2025deepseek} have demonstrated that reinforcement-learning-based post‑training can substantially enhance LLM reasoning. Motivated by this reinforcement learning with verifiable rewards (RLVR) (e.g.~\citet{shao2024deepseekmath,guo2025deepseek,lambert2024tulu,yang2025qwen3,hu2025open,yu2025dapo,guan2025rstar,zeng2025simplerl} among many other works) has become a major approach for post‑training LLMs to improve reasoning. A broad literature studies how to make RLVR training effective and efficient at scale, including novel reinforcement learning algorithms and objectives~\citep{yu2025dapo,liu2025understanding,yue2025vapo,zheng2025group}, verifier architecture and reward deigns~\citep{zuo2025ttrl,zhao2025learning,agarwal2025unreasonable,prabhudesai2025maximizing}, and mechanisms that manage exploration diversity and entropy~\citep{cheng2025reasoning,chen2025pass,cui2025entropy,wang2025beyond}, or takes a critical view on the current evaluation~\citep{yue2025does,zhao2025echo,hochlehnert2025sober}. Despite steady progress, a fundamental challenge is to balance exploration and exploitation along long reasoning trajectories without brittle heuristics. We adopt a simple yet effective annealed sampling schedule that front‑loads exploration and cools later steps during rollout to encourage exploration while keeping training stable.

\paragraph{Exploration Control in RLVR.} A line of works that is close to our work studies how to control exploration and sampling while doing RLVR~\citep{hou2025t1, cheng2025reasoning,tangoptimizing,chen2025pass,cui2025entropy,wang2025beyond,xiong2025minimalist,deng2025trial,xu2025not,zheng2025act,dou2025improving, li2025treepo}. In particular, \citet{cheng2025reasoning} propose per-token entropy to focus exploration at branching tokens in sampling; \revise{\citet{tangoptimizing, chen2025pass} transform per-prompt rewards to optimize pass@$k$}, guiding exploration across samples; \citet{cui2025entropy} control high-covariance tokens to prevent entropy collapse and sustain exploration; \citet{wang2025beyond} update only high-entropy tokens, concentrating exploration where decisions split. \revise{While more nuanced multi-threaded exploration strategies~\citep{pan2025learning} could benefit training, adapting them to the training loop for large models often introduces significant computational overhead and the challenges of maintaining massive parallelization.} Different from prior work, this paper presents the first systematic analysis of sampling temperature and introduce a purely sampling-level annealed schedule that encourages exploration and then progressively stabilize answers, thereby enabling discovery of new solutions while yielding more stable training.

% \ych{@chenxiao, perhaps also \citep{hou2025t1}, \citep{li2025treepo} and \citep{deng2025trial}? Also I am not sure we can claim ourselves as ``systematic analysis of sampling temperature''?}

\paragraph{Simulated Annealing.} 
% The idea of annealing can trace back to 
Simulated Annealing (SA) is a probabilistic optimization technique inspired by annealing in metallurgy, designed to find the global optimum in a large search space \citep{kirkpatrick1983optimization, bertsimas1993simulated}. The core principle involves a temperature parameter that controls the probability of accepting suboptimal states. Initially, a high temperature allows the search to escape local minima by exploring broadly (exploration). As the temperature gradually decreases, the algorithm increasingly favors better states, converging towards a high-quality solution (exploitation).
This ``coarse-to-fine" search dynamic, where high temperatures establish a solution's general structure and low temperatures refine its details, strongly parallels the generative process of LLMs \citep{yang2025alignment}. SA has been adapted in various machine learning contexts to manage the exploration-exploitation trade-off, including recent applications in graph optimization~\citep{liu2021simulated}, text editing \citep{zhang2024edt}, non-autoregressive generation \citep{israel2025enabling}, and efficient Best-of-N sampling \citep{manvi2024adaptive}. However, these prior applications invariably apply a single, uniform temperature across all positions in a generated sequence. This approach fails to account for the heterogeneous roles of tokens at different positions \citep{wang2025beyond}. Our work departs from this convention. To the best of our knowledge, we are the first to introduce an \textbf{intra-sequence annealed temperature} schedule, where the temperature varies dynamically within the generation of a single sequence. This novel approach allows for more nuanced control over exploration and leads to significant performance gains in RLVR.
% This simulated annealing technique is designed for naturally guiding the sampling process to explore high-entropy area first and gradually move to more deterministic low-entropy areas.